\section{Exact tests for a block of a split orthogonal matrix}
\label{sec:orthogonal}

The finite Fourier argument in
\path{grh/src/New/Spanning/SecondMoment/PositiveType.lean}
is independent of the final zero-free theorem. Its estimate
\texttt{re\_sum\_addChar\_dotProduct\_mul\_le} bounds an additive character
coefficient through the masses of two stabilizers; its final theorem
\texttt{card\_mul\_sum\_det\_eq\_zero\_le} applies this to a singular
matrix compression. Earlier ProveIt reports already passed from vectors
to frames, obtained rectangular matrix distribution bounds, and computed
the precise obstruction for isotropic blocks of symplectic matrices.
We do not count those results again. Here is a different exact boundary
calculation for split orthogonal matrices. It identifies an optimal
statistical test for the \emph{entire block}, not only its rank.

This result is a complete human proof. The group counts it uses are
classical; for related counts of isotropic subspaces see Li--Zhang,
Lemma~3.2.2, and for finite-field rank laws see Fulman--Goldstein.\cite{lizhang,fulman}
The exact block comparison below is new to the inspected report corpus.
No claim of worldwide priority is made.

\subsection{Statement}

Let $q$ be an odd prime power, $m\ge1$, and put
\[
 J=\begin{pmatrix}0&I_m\\I_m&0\end{pmatrix},\qquad
 G=\{g\in\operatorname{GL}_{2m}(\mathbb F_q):g^{\mathsf T}Jg=J\}.
\]
This is the split orthogonal group $O^+(2m,q)$.
Its determinant-one subgroup is denoted by $G^+=SO^+(2m,q)$; this
notation does not mean the possibly smaller commutator subgroup.
Write
\[
 g=\begin{pmatrix}A&B\\C&D\end{pmatrix}
\]
in this hyperbolic basis, and let $U_m$ be a uniform $m\times m$ matrix
over $\mathbb F_q$. Thus the entries of $U_m$ are independent and uniform.
For laws $P,Q$ on the finite matrix space,
$d_{\mathrm{TV}}(P,Q)=\frac12\sum_M|P(M)-Q(M)|$.

\begin{theorem}[Optimal parity and full-rank tests]
\label{thm:orthogonal-optimal-test}
If $g$ is uniform in $G^+$, then $\operatorname{rank}C$ is even and
\begin{equation}
 d_{\mathrm{TV}}(C,U_m)
   =\Pr\bigl(\operatorname{rank}U_m\text{ is odd}\bigr).
 \label{eq:SO-TV}
\end{equation}
Consequently the test ``the block has odd rank'' achieves the full
total variation distance; no test using all its entries improves this
distinguishing advantage.

If instead $g$ is uniform in the full group $G$, then
\begin{equation}
 d_{\mathrm{TV}}(C,U_m)
  =\prod_{j=1}^{m}(1-q^{-j})
    -\frac{1}{2\prod_{j=1}^{m-1}(1+q^{-j})}.
 \label{eq:O-TV}
\end{equation}
Here the full-rank test achieves the full distance. Empty products equal
one, so the formula also includes $m=1$.
\end{theorem}

For clarity, ``achieves the full distance'' uses
$d_{\mathrm{TV}}(P,Q)=\max_E|P(E)-Q(E)|$.
With equally likely hypotheses, the corresponding best success
probability is $(1+d_{\mathrm{TV}})/2$.

\subsection{Proof by counting isotropic subspaces}

Let $W=\mathbb F_q^m\oplus0$ and $L=gW$.
The space $L$ is maximal totally isotropic for $J$ and
\[
 \operatorname{rank}C=m-\dim(L\cap W).
\]
Write
\[
 {m\brack r}_q
 =\prod_{i=0}^{r-1}\frac{q^m-q^i}{q^r-q^i},\quad
 F_q(m,r)=\prod_{i=0}^{r-1}(q^m-q^i),\quad
 T_m=2\prod_{i=1}^{m-1}(1+q^i).
\]

\paragraph{The determinant and rank parity.}
Every element of $G$ stabilizing $W$ has block form
\(
 \left(\begin{smallmatrix}P&*\\0&P^{-\mathsf T}\end{smallmatrix}\right)
\), hence determinant one.
Suppose $\dim(L\cap W)=m-r$. A change of basis in $W$, extended by the
inverse transpose on its complementary isotropic subspace, puts the
intersection in standard position. In its orthogonal quotient, $L$ is
transverse to the remaining part of $W$, hence is the graph of a
skew-symmetric matrix. An upper unipotent isometry with a skew-symmetric
upper-right block removes that graph. Both changes stabilize $W$ and
have determinant one. The resulting space is
\[
 \langle e_1,\ldots,e_{m-r},f_{m-r+1},\ldots,f_m\rangle.
\]
It is the image of $W$ under the isometry which swaps $e_i$ with $f_i$
in precisely $r$ hyperbolic pairs. That isometry has determinant
$(-1)^r$. Multiplying by it leaves an element stabilizing $W$, proving
\begin{equation}
 \det g=(-1)^{\operatorname{rank}C}.
 \label{eq:orthogonal-parity}
\end{equation}
This construction also proves transitivity on maximal isotropic
subspaces. Since their stabilizer has determinant one, the two determinant
cosets of $G$ give two equally sized families, selected by the parity
of $r$.

\paragraph{The exact rank counts.}
Choose $D=L\cap W$ of dimension $m-r$ in ${m\brack r}_q$ ways.
The quotient $D^\perp/D$ is a hyperbolic space of dimension $2r$.
The possible $L/D$ are precisely its maximal isotropic complements to
$W/D$. Writing each such complement as a graph gives a skew-symmetric
$r\times r$ matrix. There are $q^{r(r-1)/2}$ such matrices because
$q$ is odd. Thus the number of possible $L$ at block rank $r$ is
\begin{equation}
 N_m(r)={m\brack r}_q q^{r(r-1)/2}.
 \label{eq:orthogonal-subspace-count}
\end{equation}
The elementary $q$-binomial identity
\[
 \sum_{r=0}^{m}{m\brack r}_q q^{r(r-1)/2}
   =\prod_{i=0}^{m-1}(1+q^i)=T_m
\]
follows by induction from the Gaussian-binomial recursion.
Uniform group elements induce the uniform law on each orbit. Therefore
the full orthogonal rank law and the special orthogonal rank law are
\begin{align}
 o_m(r)&=N_m(r)/T_m,\label{eq:O-rank-law}\\
 s_m(r)&=\begin{cases}2o_m(r),&r\text{ even},\\0,&r\text{ odd}.
                 \end{cases}\label{eq:SO-rank-law}
\end{align}
For an independent-entry matrix, choose its $r$-dimensional image and
then a surjection onto it. This gives
\begin{equation}
 u_m(r)=\Pr(\operatorname{rank}U_m=r)
       ={m\brack r}_qF_q(m,r)q^{-m^2}.
 \label{eq:iid-rank-law}
\end{equation}

\paragraph{Why rank suffices for all matrix tests.}
For each $P\in\operatorname{GL}_m(\mathbb F_q)$,
the matrix $\operatorname{diag}(P,P^{-\mathsf T})$ lies in $G^+$.
Left and right multiplication by these matrices transform $C$ into
$P^{-\mathsf T}CQ$ for independently arbitrary invertible $P,Q$.
Uniform measure on either group is invariant under these operations.
Since left and right invertible transformations act transitively on
matrices of a given rank, both block laws are uniform conditional on
rank. The independent-entry law has the same property. Thus the total
variation distance of the full matrices equals that of their rank laws.

\paragraph{Comparing the mass at each rank.}
Set $R_m(r)=2o_m(r)/u_m(r)$, also for odd $r$. Cancellation in the exact
formulas gives
\[
 R_m(m)=\frac{1}{(1-q^{-m})
                      \prod_{j=1}^{m-1}(1-q^{-2j})}>1,
 \qquad
 \frac{R_m(r-1)}{R_m(r)}=q^{m-r+1}-1.
\]
Because $q\ge3$, each displayed ratio is at least two.
Hence $s_m(r)>u_m(r)$ at every even rank, while $s_m(r)=0$ at every
odd rank. All the deficit lies at the odd ranks, proving
\eqref{eq:SO-TV}.

For the full orthogonal group, the likelihood ratio is $R_m(r)/2$.
At rank $m$ it is less than one: writing $Q=q^{-1}\le1/3$, its
denominator before the factor two satisfies
\[
 (1-Q^m)\prod_{j=1}^{m-1}(1-Q^{2j})
 \ge (1-Q)\left(1-\sum_{j=1}^{\infty}Q^{2j}\right)
 \ge \frac{7}{12}>\frac12.
\]
It is greater than $1/2$, because that same product is less than one.
At rank $m-1$ the ratio is therefore greater than
$(q-1)/2\ge1$, and it increases as rank decreases.
Thus only full rank has deficient mass in $o_m$ compared with $u_m$.
Their full-rank probabilities are respectively
\[
 o_m(m)=\frac{1}{2\prod_{j=1}^{m-1}(1+q^{-j})},\qquad
 u_m(m)=\prod_{j=1}^{m}(1-q^{-j}).
\]
Their difference proves~\eqref{eq:O-TV}.\hfill$\square$

\subsection{Finite examples and the limiting obstruction}

The supporting program \path{code/check_orthogonal_blocks.py} uses exact rational
arithmetic. It independently enumerates all maximal isotropic subspaces
and all square matrices for $q=3$, $m=1,2,3$. At $m=3$ it examines all
$33{,}880$ three-dimensional subspaces of $\mathbb F_3^6$ in unique
reduced-row-echelon form, of which $80$ are isotropic. It also independently
enumerates all $2$ and $576$ elements of $SO^+(2m,3)$ for $m=1,2$ from
elementary generators, checks the preserved form, and checks every
individual block probability. No group enumeration at $m=3$ is claimed.

\begin{center}
\begin{tabular}{c|c|c|l}
 $m$ & $d_{\mathrm{TV}}(C_{SO},U_m)$ &
       $d_{\mathrm{TV}}(C_O,U_m)$ & nonzero $SO$ rank masses\\\hline
 $1$ & $2/3$ & $1/6$ & $r=0:\ 1$\\
 $2$ & $32/81$ & $47/216$ & $r=0:\ 1/4,\quad r=2:\ 3/4$\\
 $3$ & $11570/19683$ & $13597/58320$ &
       $r=0:\ 1/40,\quad r=2:\ 39/40$
\end{tabular}
\end{center}

For the full orthogonal group, fixed $q$ and $m\to\infty$ in
\eqref{eq:O-TV} give the positive limiting distance
\[
 \prod_{j=1}^{\infty}(1-q^{-j})
 -\frac{1}{2\prod_{j=1}^{\infty}(1+q^{-j})}.
\]
The products converge to positive limits, and the proof above shows
that the difference is positive for every odd $q$.
Even when $q\to\infty$, this expression tends to $1/2$.
The matrix entries therefore do not become independent at this maximal
isotropic boundary. The determinant component leaves a precise parity
signature, and averaging the two components leaves a precise full-rank
signature. This is a method transfer in finite geometry and probability,
with no use of RH or GRH as a premise.

